\documentstyle[psfig,12pt]{article}
\input define
%%%%%%%%%%%%%%%%%%%%%%%%%%%%%%%%%%%%%%%%%%%%%%%%%%%%%%%%%%%%%%%%%%%%%%%%%%%%
%                                                                          %
%       Numerical differentiation in system theory. Observer design        %
%                                                                          %
%       Edited on August 8, 1993                                           %
%	Converted by Laura to LaTex on Sept. 6, 1993                       %
%       Section on Examples added on Sept. 7, 1993                         %
%%%%%%%%%%%%%%%%%%%%%%%%%%%%%%%%%%%%%%%%%%%%%%%%%%%%%%%%%%%%%%%%%%%%%%%%%%%%
%
%%%%%%%%%%%%%%%%%%%%%%%%%%%%%%%%%%%%%%%%%%%%%%%%%%%%%%%%%%%%%%%%%%%%%%%%%%%%
%   Comment out and adjust the following lines when using Double Columns   %
%%%%%%%%%%%%%%%%%%%%%%%%%%%%%%%%%%%%%%%%%%%%%%%%%%%%%%%%%%%%%%%%%%%%%%%%%%%%
% \hoffset = -10 true mm
% \voffset = -5 true mm
% \hsize = 185 true mm
% \vsize = 235 true mm
% \newdimen\separatorwidth
% \separatorwidth = 5 true mm
% \tolerance = 1000
% \emergencystretch = 25pt
%%%%%%%%%%%%%%%%%%%%%%%%%%%%%%%%%%%%%%%%%%%%%%%%%%%%%%%%%%%%%%%%%%%%%%%%%%%%
%
%\documentstyle[proc,acc]{article}
%\pagestyle{empty}
%%%%%%%%%%%%%%%%%%%%%%%%%%%%%%%%%%%%%%%%%%%%%%%%%%%%%%%%%%%%%%%%%%%%%%%%%%%
%                                                                         %
%	Commented lines below are for two columns                         %
%                                                                         %
%%%%%%%%%%%%%%%%%%%%%%%%%%%%%%%%%%%%%%%%%%%%%%%%%%%%%%%%%%%%%%%%%%%%%%%%%%%
%\textheight=24cm
%\textwidth=17.7cm
%\oddsidemargin=-.4cm
%\columnsep=1.5cm
%\topmargin=-2.716cm
\textheight=9in
\textwidth=7.0in
\oddsidemargin=0in
\topmargin=-.66in

\def\thesection{\arabic{section}.}
\def\thesubsection{\thesection\arabic{subsection}.}
\def\theequation{\arabic{equation}}
\def\figurename{\small Figure}

\newcommand{\r}{\hbox{\sl l\kern-.2em R}}
\def\thesection{\arabic{section}.}
\def\theequation{\thesection\arabic{equation}}
\def\theenumi{\roman{enumi}}\def\labelenumi{(\theenumi).}
\def\theenumii{\alph{enumii}}\def\labelenumii{(\theenumii)}
\newtheorem{theorem}{Theorem}[section]
\newtheorem{proof}{Proof:}
\def\theproof{}
\def\thetheorem{\thesection\arabic{theorem}}

\def\setR{\hbox{$l\! \! R$}}
\def\setN{\hbox{$l\! \! N$}}
\def\CN{{\cal N}}

\def \BA {{\bf A}}
\def \BB {{\bf B}}
\def \BC {{\bf C}}
\def \BE {{\bf E}}
\def \BF {{\bf F}}
\def \BG {{\bf G}}
\def \BH {{\bf H}}
\def \BI {{\bf I}}
\def \BJ {{\bf J}}
\def \Bk {{\bf k}}
\def \BK {{\bf K}}
\def \BL {{\bf L}}
\def \BM {{\bf M}}
\def \BN {{\bf N}}
\def \BQ {{\bf Q}}
\def \BP {{\bf P}}
\def \BQ {{\bf Q}}
\def \BR {{\bf R}}
\def \BS {{\bf S}}
\def \BV {{\bf V}}

%\def \CN   {{\eusm N}}

\def\thesubsubsection{\thesubsection\arabic{subsubsection}.}

\begin{document}

\title{Interpolation and Numerical Differentiation \\
for Observer Design}

\author{S. Diop\thanks{Charg\'e de Recherche CNRS - LAGEP, UCB Lyon I - 43 Bd.
11 Nov. 1918, 69622 Villeurbanne France;
work supported in part by a grant from NSF/CNRS and in part 
by the Control Systems Laboratory of
the Dept. Elect. Eng. and Comp. Sci., U. of Michigan.}, 
J.\ W.\ Grizzle\thanks{Control Systems Laboratory, \
Department of Electrical Engineering and Computer
Science, \ University of Michigan, Ann Arbor, MI 48109-2122;
work supported in part by the National Science
Foundation under contract NSF ECS-92-13551; matching funds to this grant were
provided by the FORD MO.\ CO.}
P.\ E.\ Moraal\footnotemark[2] and A. Stefanopoulou\footnotemark[2] }
%\\[-48pt]}

\maketitle
%\def\thefootnote{\fnsymbol{footnote}}
%\footnotetext[1]{\mbox{}\footnotemark[1] Laboratoire d'Automatique et de G\'{e}nie des
%Proc\'{e}d\'{e}s, CNRS-Universit\'{e} Claude Bernard, B\^{a}t. 721,
%43 Bd. du 11 Novembre 1918, 69622 Villeurbane-France; work supported in part by
%a NATO Postdoctoral Fellowship. }

%\footnotetext[2]{\footnotemark[2] Department of Electrical Engineering and Computer
%Science,University of Michigan, Ann Arbor, MI 48109-2122;
%work supported in part by the National Science
%Foundation under contract NSF ECS-88-96136; matching funds to this grant were
%provided by the FORD MO.\ CO.}


\abstract  
This paper explores interpolation and numerical differentiation
as a basis for constructing a new approach to the design of nonlinear
observers.  
Numerical differentiation is well known to be an ill-conditioned problem in the
sense that small perturbations on the function to be differentiated may induce
large changes in the derivatives.
We illustrate through examples that some specific methods may be able to overcome this
difficulty, turning numerical differentiation techniques into a possible
challenger to existing observer design methods.

\medskip\noindent
{\bf Keywords:}
Observer design; Interpolation; Numerical differentiation; System theory. 

\section{Introduction.}  Observer design is a fundamental problem in system theory and 
control practice.  Indeed, many feedback  design techniques, especially for nonlinear
system models,  start with the assumption that the full state is available for
feedback and that the {\it user} will design separately a suitable observer
for the actual implementation.  This in turn has motivated a great deal of
research into observer design.   

The specific problem addressed herein involves the situation where
one seeks to ``estimate'' the states of a continuous-time model from
observations collected at discrete instants in time.  To date, this problem has
been addressed only in the context of sampled data systems
\cite{grizzle1990b,grizzle1990a,grizzle1990c,moraal1992}, where the
procedure has been to first compute
a discrete-time model from the continuous-time model, and
then to design an observer on the basis of the discretized system model.
Of course, the discretization process sometimes can be quite non-trivial, from
a numerical or computational point of view; also, it is usually quite important 
to assume that the observed data has been collected at regular intervals of time,
and thus dealing with problems where event driven sampling is necessary can be difficult.  

This paper comes out of the observation that
on the one hand, the numerical analysis literature contains a
substantial body of results on numerical differentiation, while on
the other hand, one of the very basic problems in system theory turns
out to be that of obtaining the time derivatives of a continuous system
variable, which is usually known only through the differential equations of
the system, and time samples of this variable.
We present a first attempt in applying results on numerical differentiation
to the observer design problem, mainly by
working out some specific examples.

To see how numerical differentiation and observer design may tie together, 
let us consider the 
following simple observer design problem.
Let a system be given as
\begin{equation}
\left\{
\begin{array}{rcl}
\dot x_1 & = & x_2\\
\dot x_2 & = & -ux_1^3\\
y & = & x_1\end{array}
\right. \label{1}
\end{equation}
and let the question be to build an observer for (\ref{1}); that is,
to reconstruct the states of the system on the basis of the available
measurements.
It is clear that the variable $x_1$ need not be computed since it 
is directly given by the measurement: $ x_1 = y$.
From the equations of (\ref{1}), we see that $x_2 = \dot y$.
If, as is usually the case, we have no access to 
$\dot y$ then we need to {\it compute} $\dot y$ from the measurement $y$.
This is where numerical differentiation comes into play:
we explore the feasibility of {\it numerically
differentiating} $y$ in order to obtain $x_2$.
If $y$ is known only through its time samples then numerical 
differentiation yields an approximation of
$\dot y$ at the current sampling time $t_k$ by some number
$\hat {\dot y}(t_k)$.
An algorithm for the computation of $\hat {\dot y}(t_k)$ from the
samples of $y$, and {\it estimates of the error bounds} should thus be worked out.
Usually, not only $\dot y$, but, also, finitely many higher
derivatives of $y$ should be approximated at time $t_k$ in order to be
able to compute all of the desired state components.

The intent of this simple example is to 
illustrate our main objectives in applying 
numerical differentiation techniques to the observer design problem.
If efficient numerical 
differentiation algorithms are available, then they should be
applicable to many other system theoretic problems that require the
computation of the derivatives of continuous variables which are
known only through their time samples and the differential equations they
satisfy.
The inversion problem is among other potential applications that come to mind.

Our investigation led us directly to interpolation theory, and more
generally, to approximation theory.
As is well known, there is a considerable number of different ways for
approximating the value of an unknown function at some point.
One commonly chooses among polynomial, rational, periodic, or exponential
functions
as basis elements.
(See also methods stemming from the Shannon sampling theorem
\cite{butzer1971,butzer1977,splettstosser1979,splettstosser1983,butzer1983a,butzer1983b,ries1983,schussler1983}.)
The optimality of any approximation process is, of course, relative to the
assumed basis functions.
Approximation methods based on polynomials are the
simplest ones, but spline approximations seem to have some interesting
fundamental properties.


%{\it Tornambe: talk about his work.}
%As far as we know, system theorists used to invoke asymptotic 
%observer design techniques to solve this problem.
%While this approach is very successful for constant linear systems, the
%nonlinear case seems restrictive, and even obstructive.

The paper is organized as follows.  Section \ref{observability} briefly reviews
the notion of observability we are using.  Section \ref{numdiff} provides a 
considerable review of the numerical analysis literature on numerical
differentiation.  Section \ref{examples} illustrates on specific examples 
how numerical
differentiation
may be integrated into an overall observer design strategy.  
%A few concluding remarks are given in Section \ref{conclusions}.


\section{An observability property.}\setcounter{equation}{0}
\label{observability}

Given a system described by the state equations
\begin{equation}
\left\{
\begin{array}{rcl}
\dot x_i & = g_i(x, u) \quad (1 \le i \le n)\\
y_j & = h_j(x, u) \quad (1 \le j \le p)\end{array}
\right.
\label{2}
\end{equation}
we will adopt the following specific observability property (see
\cite{chung1990,grizzle1990b,grizzle1990c,grizzle1990a,moraal1992,song1992,moraal1993b,moraal1993a}):
system (\ref{2}) is {\bf observable} if there is an integer $N$ such
that the map $H\left(u, \dot u, \cdots, u^{(N-2)}, \cdot \right)$ of
the state space into some space of output values, defined by
\begin{equation}
H\left(u, \dot u, \cdots, u^{(N-1)}, x \right) = 
\left( 
\begin{array}{c} y \\ \dot
y \\ \vdots \\ y^{(N-1)} \end{array}\right), \label{3}
\end{equation}
is {\it injective} for any fixed {\it universal input}.
%We shall define later what these universal inputs are.
It is easiest at this stage to assume that there is actually
no input in (\ref{2}), or that all inputs are universal in the sense that
the above map is injective for any input. 

For an observable system (\ref{2}), we then may write
$$x = L\left(u, \dot u, \cdots, u^{(N-2)}, y, \dot y, \cdots, y^{(N-1)} 
\right).$$
The existence of the map $L$ is guaranteed by our definition of 
observability, but an explicit expression of $L$ may not be easy to 
obtain so that numerical equation solving methods such as Newton's algorithm
may be required in order to obtain $L$.  
When (\ref{2}) is a rational system in the sense that the $g_i$'s and the 
$h_j$'s are rational functions of their arguments, then the type of 
observability we just defined corresponds to {\it rational 
observability} of \cite{diop1991d,diop1992b}, and then the explicit
expression of $L$ may be  constructively derived through the use of certain
differential
algebraic techniques.

%\noindent {\bf Note:} In the literature, it is often tacitly assumed that $N$ should
%be taken  as $2n + 1$; we shall examine this situation which, {\it a
%priori}, seems connected to the works \cite{aeyels1981a,aeyels1981b}.

For these observable systems, the observer design problem may thus be 
seen as a problem of numerical differentiation: once estimates of the
derivatives
of $y$ and $u$ can be determined from the available measurements, then $x$ can be
determined from $L$.

\section{On numerical differentiation.}\setcounter{equation}{0}
\label{numdiff}

The basic problem we are facing may be roughly stated as follows:

\noindent {\bf Problem P.}
A real valued function $f$ of one real variable
(called the  {\it time,} $t$) defined on an interval $\BI$ is known
only through the differential equations it satisfies and through the
values it  takes at some instants $t_0, t_1, \cdots, t_k \in \BI$,
with $t_0 <$ $t_1 < \cdots <$ $t_k$.  The question is to provide
estimates of the first  $q$ derivatives $f^{(1)}(t_k)$,
$f^{(2)}(t_k), \cdots$, $f^{(q)}(t_k)$ of $f$ at time $t_j$, for some $0 \le j
\le k$.
The difficulty of the problem may be increased by the presence of noise in the samples
of $f$, in which case, the approximation technique should provide
{\it smoothing}, and even the values of $f$ would need to be estimated.

\noindent {\bf Remarks:}

\begin{itemize}

\item
The more general problem we should end up facing in system
theory is actually a multivariable one in that $f$ is a vector valued function
instead of scalar valued as in the latter statement.

\item
If $r$ is the order of the differential equations defining $f$ then only
at most the first $q = r - 1$ derivatives of $f$ need be estimated, the
succeeding derivatives being computable from the previous ones
by means of these differential equations defining $f$.

\item
Problem~P is a particular case of the general problem of approximation of
functions on the basis of partial knowledge of these functions.
Its seems that there is not yet a complete solution with
expressions of the error bounds in the literature.
One partial solution consists in taking the derivatives of the
{\it interpolant} $\hat f$ for those of $f$.
On this approach, there is a substantial amount of literature.
Another possible direction consists in trying to overcome
the ill-conditionedness of the numerical differentiation problem by using
some regularization techniques.
The only known works are the one by Cullum~\cite{cullum1971},
and those in the Russian literature (cited in the latter paper) to which we
do not yet have access.
\end{itemize}

\subsection{Generalities on approximation theory.}
The basic problem that approximation theory addresses is the following:
Assume that the unknown object to be approximated is an element
$x$ of a given metric space $\BE$, and find a {\it best estimate}
$\hat x$ of $x$ in an (also given) subspace $\BF$ of $\BE$ subject to
$$
{\rm d}(x,\hat x) = {\rm d}(x,\BF) = \min_{y \in \BF}{\rm d}(x,y),
$$
where {\rm d} denotes the distance on $\BE$.
The word {\it best} is of course relative to the distance used in
$\BE$ and to the choice of $\BF$.
For example, for a fixed distance, the larger $\BF$ is, the finer the
approximation will be.
This problem has at least one solution if the subspace
$\BF$ is compact, or if $\BE$ is a normed vector
space and $\BF$ is a linear subspace of finite dimension, or if $\BE$
is a Hilbert space and $\BF$ is a closed subspace.

In our case, the object to be approximated is the set of derivatives $f^{(1)}$,
$f^{(2)}, \cdots$, $f^{(q)}$ of a real valued function $f$ defined
on, say a compact interval $\BI$ of $\setR$, so that the space $\BE$
may be taken as the set of real valued functions which are $q-1$ times
differentiable with $f^{(q-1)}$ absolutely continuous and $f^{(q)}$ square
integrable on $\BI$, or which are infinitely differentiable, or, even, which
are analytic on $\BI$.
$\BE$ is thus endowed with its usual $\setR$-algebra
structure and may be converted into a real normed vector space by one of
the following 
\begin{eqnarray*}
\| g \|_p & = & \left( \int_\BI |g(t)|^pdt \right)^{1/p}, \, g \in \BE, p
\in \setR, \, p \ge 1, \\
\| g \|_\infty & = &\sup_{t \in \BI} |g(t)|, \, g \in \BE, \\
\| g \|_{(q)} & = &\| g \|_2 + \| g^{(1)} \|_2 + \cdots +
\| g^{(q)}\|_2, \, g \in \BE, \\
& \hbox{\rm etc.} &
\end{eqnarray*}
The norm $\| \, \|_2$ corresponds to the Hilbert space
structure on $\BE$, the corresponding inner product being
$$
\left( g, h \right) = \int_I g(t)h(t)dt, \, g, h \in \BE.
$$
The norm $\| \, \|_\infty$ is usually called the {\it infinity} or {\it
uniform} norm.

Whenever $\BF$ is a finite dimensional subspace of
$\BE$ and the norm
$\| \, \|_2$ is used,
 the {\it least squares} approximation process results in the existence and 
uniqueness of the best
estimate of any given $f$ in $\BE$.

The subset $\BF$ of $\BE$ is usually taken as a subspace of one of
the following linear subspaces:
\begin{itemize}
\item
$\setR\left[ t \right]$, the subspace of polynomial functions,
\item
$\setR\left( t \right)$, the subspace of rational functions,
\item
$\setR\left[ \cos t, \sin t \right]$, the subspace of $\setR$-linear
combinations of $\cos(it)$ and $\sin(it)$, $i \in \setN$.
\end{itemize}
%Assume linear, uniform, approximation theory,
Assume that $\BF$ is finite dimensional, with dimension $N'+1$, and norm $\| \ \|$.
We then know that, given any $f_0$ in $\BE$, there is at least one
best estimate $\hat f$ of $f_0$ in the sense of minimizing $\| f-f_0 \|$ over
$f \in \BF$.  When  $\| \ \| = \| \, \|_\infty$, $\hat f$ is called a best uniform
estimate.
A base $(\zeta_0$, $\zeta_1, \ldots$, $\zeta_{N'})$ of $\BF$ is said to
satisfy the Haar condition if any real linear combination
$$
c_1\zeta_1(t) + \cdots + c_N\zeta_{N'}(t)
$$
vanishes at less than $N'+1$ distinct points in $\BI$.

\begin{theorem}[\rm (Haar), \cite{delavalleepoussin1970}].\label{thm1}
A necessary and sufficient condition for the best uniform estimate to be
unique is that the base $(\zeta_0$, $\zeta_1, \ldots$, $\zeta_{N'})$ of $\BF$
satisfies the Haar condition.
\end{theorem}

\subsection{Differentiating the interpolant.}
Since the interpolation problem has been given a large number of elegant
and practical solutions, the easiest way to approach Problem~P is to take
the derivatives of the interpolant $\hat f$ for those of $f$.
Let us make the following remark at this stage even if we have not yet
introduced the material which makes it clear:
There might be, due the reality of a particular
physical problem, no freedom in the choice of the sampling instants $t_i$,
but the error bounds for the estimated derivatives are
reasonably expected to depend on both the number of
samples and  their distribution on the time interval $\BI$ as apparent from
the Tchebychev theorem.
Our study should address this issue, as well, but
for now, we shall suppose that the current time is $t_N$; later on, we shall
discuss how we should deal with the succeeding sampling instants $t_{N+1},
\ldots$ as time evolves.

Let $\BF_{N'}$ be a linear subspace of $\BE$ of finite dimension $N'+1$, and
with base $(\zeta_0$, $\zeta_1, \ldots$, $\zeta_{N'})$.
An element
$$
\hat \zeta_{N'}(t) = c_0\zeta_0(t) + \cdots + c_{N'}\zeta_{N'}(t)
$$
of $\BF_{N'}$ {\it interpolates} $f$ if
$$
\hat \zeta_{N'}(t_i) = f(t_i) \, (0 \le i \le N),
$$
that is, if
$$
\sum_{j=1}^{N'}c_j\zeta_j(t_i) = f(t_i) \, (0 \le i \le N).
$$
Therefore,

\begin{theorem}\label{thm2}
There is one and only one function in $\BF_{N'}$ interpolating $f$
if and only if the previous set of equations in the $c_j$'s has one and only one
solution, that is, the existence and uniqueness in $\BF_{N'}$ of an
interpolant of $f$ is equivalent to $N' = N$ and the determinant of the
following matrix
$$\left( \begin{array}{ccc}
\zeta_0(t_0)&\ldots&\zeta_N(t_0)\\
         \vdots&&\vdots\\
         \zeta_0(t_N)&\ldots&\zeta_N(t_N)\end{array} \right)$$
is nonsingular; in this case, we may use the well-known Cramer formula to
compute the interpolant, and write it in the form
$$
\hat \zeta_N(t) = \sum_{i=1}^{N}f(t_i)s_i(t) \, (t \in \BI),
$$
with $s_i(t_j) = \delta_{ij}$, where $\delta_{ij}$ is the Kronecker symbol.
\end{theorem}

\noindent
Note that the Haar condition guarantees that the latter square matrix is
nonsingular.
Note also that this previous result shows that, when interpolating linearly a
function at $N$ points, the subspace of projection, $\BF$, should be of
dimension precisely $N$.

Among the well-established results on interpolation theory, we will now discuss
those due to Lagrange, those based on splines, etc.


\subsubsection{Differentiating the Lagrange interpolant.}

It is clear that there are infinitely many polynomials of arbitrary degree
{\it interpolating} $f$, but, according to the previous general result,
there is one and only one of degree at most $N$ which interpolates $f$.
This unique polynomial $L_N$ is known as the Lagrange (or polynomial)
interpolant of $f$.

Assuming
$$
\ell_i(t) = \prod_{\scriptstyle j = 0 \atop \scriptstyle j \ne i}^N{t-t_j \over
t_i - t_j},
$$
we have
$$
L_N(t) = \sum_{i=0}^{N}f(t_i)\ell_i(t) \, (t \in \BI),
$$
with $\ell_i(t_j) = \delta_{ij}$.
This is not the most interesting form for the computation of the Lagrange
interpolant.
The Newton formulae lead to more efficient computations,
see~\cite{stoer1980,marchuk1982,nonweiler1984} for example.

If the only knowledge we have on $f$ is the set of its samples, then being
able to estimate the error bounds is hopeless. Given more information,
the following result is known:
If $f$ is $N$ times continuously differentiable in $\BI$, then for any $t$
in $\BI$, there is a point $\xi_t$ in $\left[ t_0, t_N \right]$ such that
$$
\varepsilon(t) = f(t) - \hat \zeta_{N}(t) = {1 \over
(N+1)!}\ell(t)f^{(N+1)}(\xi_t) \, (t \in \BI),
$$
where $\ell(t) = \prod_{i=0}^N(t-t_i)$.

This follows mainly from the Rolle theorem.
In general, given $t$, we do not know $\xi_t$ nor
$f^{(N+1)}(\xi_t)$.
This is the reason why we shall need to know a bound for the higher
derivatives of $f$:
$$
\| \varepsilon \|_\infty \le {\Delta^N \over (N+1)!}\| f^{(N+1)} \|_\infty,
$$
where $\Delta$ is $\max_{1 \le i \le N}(t_i -t_{i-1})$.
This bound depends both on the number $N+1$ of samples and on their
distribution.
If it is possible to choose freely the sampling instants, then we would be
able to decrease the latter error bound by sampling according to the
Tchebychev theorem, see \cite{nonweiler1984}.
Now, the error when taking $\hat \zeta_{N}^{(j)}(t)$ for $f^{(j)}(t)$ is
given by
$$
\varepsilon^{(j)}(t) = f^{(j)}(t) - \hat \zeta_{N}^{(j)}(t) = \sum_{\iota =
0}^j{1 \over (N+\iota+1)!}{j \choose
\iota}\ell^{(j-\iota)}(t)f^{(N+\iota+1)}(\xi_\iota) \, (t \in \BI),
$$
where the $\xi_\iota$'s are $t$ dependent and in $\left[ t_0, t_N
\right]$.
Therefore, we have
$$
\| \varepsilon^{(j)} \|_\infty \le c \, \Delta^{N+1-j} \max_{0 \le \iota \le
j}\| f^{(N+1+\iota)} \|_\infty,
$$
where $c$ is an absolute constant (it does not depend on $f$ nor
the samples), see \cite{schaback1992}.

\noindent {\bf Remark:}
Though simple and elegant, Lagrangian interpolation suffers from the
{\rm Runge phenomenon} which is the fact that a small change to the data
points in the middle of the data set may produce a very large excursion in
the end points, $t_0, t_N$, see \S 3.2.3 of \cite{nonweiler1984}.
According to the latter book,
Lagrangian interpolation is thus not recommended for large set of equispaced
data.
Nevertheless, there are some improvements of the Lagrangian interpolation in
the literature: \cite{schaback1992} addresses the problem of numerical
differentiation with application to computer-aided-design (the basic
assumption that we know the range of $f$ made in that paper makes it
unrealistic in control theory, but still deserves some attention),
\cite{criscuolo1990,criscuolo1991,criscuolo1992a,criscuolo1992b} propose an
{\it osculatory} alternative by assuming available some extra data points
which represent values of the higher derivatives of $f$ at the sampling
instants.

\subsubsection{Differentiating the spline interpolant.}

Excellent introductions to spline theory may be found in
\cite{greville1969b,nonweiler1984}; \S 3 of the latter book is particularly
recommended for a motivated step by step introduction of splines in the
theory of interpolation.

A {\it natural spline} of {\it degree} $2l-1, \, l \in \setN, \, l \ge 1$
with {\it knots} $t_0, t_1, \cdots, t_N$ is a real valued function, $s$,
defined on $\setR$, and such that:
\begin{enumerate}
\def\theenumi{\roman{enumi}}
\def\labelenumi{(\theenumi)}
\item $s$ has $2l-2$ first derivatives defined and continuous on
$\setR$,
\item In each interval $\left[ t_i,t_{i+1} \right], \, i = 0, 1,
\cdots, N-1$, $s$ is polynomial and with degree at most $2l-1$,
\item In each interval $\left] \infty,t_0 \right], \left[t_N,\infty
\right[$ $s$ is polynomial with degree at most $l-1$.
\end{enumerate}

The set of natural splines of degree $2l -1$ with knots $t_0$, $t_1,
\cdots$, $t_N$ is classically denoted by $\CN_{2l-1}(t_0,$ $t_1, \ldots$,
$t_N)$.
It is clear that $\CN_{2l-1}(t_0, t_1, \ldots, t_N)$ contains the set
of real coefficient polynomials in $t$ with degree less than $l$.

The fine points of spline interpolation theory are the following
fundamental properties, see
\cite{schoenberg1964d,deboor1966,greville1969b}, for example, for more
properties and proofs.

\noindent {\bf Interpolation property.}
Let $f$ be in $\BE$.
For any natural integer $l$ such that $l \le N+1$ there is
one and only one natural spline $\hat s$ in $\CN_{2l-1}(t_0, t_1, \ldots
t_N)$ which interpolates $f$.
When $l = N+1$ then, of course, $\hat s$ is the Lagrange polynomial which is
the unique polynomial of degree less than $N+1$ interpolating $f$.
For $l \ge N+1$, there are infinitely many natural splines which
interpolate $f$.

\noindent {\bf Smoothing properties.}
Given $f$ in $\BE$, and $l \le N+1$ then
$$
\int_\BI \left( s^{(l)}(t) \right)^2 dt \le
\int_\BI \left( f^{(l)}(t) \right)^2 dt
$$
for any $s$ in $\CN_{2l-1}(t_0, t_1, \ldots, t_N)$, and equality holds
only if $f$ is the unique natural spline $\hat s$ interpolating $f$.

The reason why this last property is called {\it smoothing} is
explained very clearly in \S 3.4 of \cite{nonweiler1984}. 
It refers to {\it cubic splines} where the above integrals may be
interpreted in terms of energy.
Note that this minimum property may be generalized a bit farther as is done
in Sard's linear approximation theory~\cite{sard1967}, see also
\cite{greville1969b}, and
\cite{micchelli1976c,micchelli1976a,micchelli1976b}.
Roughly speaking, the generalization consists in substituting a general
linear functional for $f^{(l)}(t)$ in the above integrals.
We need the theory of Peano kernels, or more generally, of reproducing
kernels in Hilbert spaces in order to understand and use these
generalizations of spline interpolation.
%Anyhow, the idea of smoothing
%should not remain, instead, we have that of minimization.
One may ask if an ultimate generalization to the nonlinear case would
not be a definite solution to Problem~P;see \cite{micchelli1976b} and
the references therein, for what
seems to be a step towards this idea. 

We postpone to the next draft of this paper the summary of the results
on error bounds of spline based, numerical differentiation.
Let us just mention some papers containing partial results:
\cite{golomb1959,varma0000,birkhoff1964,secrest1965,sard1967,greville1969b,dolezal1982}.
We have not yet consulted \cite{deboor1978b} which is reported to
be an excellent source for spline interpolation.
This last paper contains an
innovative idea which consists in improving the spline interpolation by
iterating the basic process of spline interpolation as many times as desired. 
It is claimed that arbitrary precision may
thus be obtained.

Spline interpolation certainly provides better estimation, and constitutes
the most elaborate method of interpolation when one restricts oneself to
linear approximation with polynomial base functions.
Beyond that method, one has to enter the theory of {\it nonlinear
approximation} (see~\cite{rice1965}) with its simplest form being {\it
rational approximation.}
But then, caution is recommended in the literature,
and, anyway, results are far less simple and clear.

The error bounds, say $\varepsilon_0$, $\varepsilon_1, \ldots$, on the
respective derivatives $f^{(0)}$, $f^{(1)}, \ldots$ of a function $f$,
no matter what the interpolation process is, are reasonably
expected to be such that $\varepsilon_0 <$ $\varepsilon_1 < \ldots$,
see~\cite{czipszer1958}.
%This means that the lower the derivatives of $f$ to be estimated are the
%easier we could afford high accuracy on the estimates of the derivatives.

Finally, let us note that we do yet not have at hand the references
\cite{newman0000,anderssen1972} which seem to be very
interesting\footnote{It is remarkable that these
researchers were with engineering departments such as the IBM
Research Center at Yorktown Heights, NY; some other contributors on
numerical differentiation used to be with Ford Mo.~Co., Dearborn, MI.}.

\subsection{Numerical differentiation by regularization.}
As we said earlier, the only reference on this approach that we yet have at
hand is the paper by Cullum~\cite{cullum1971}.
This paper mentions some other works by Russian numerical analysts.
Further bibliographic research is also due since there might be other
works which have been completed since 1971.

%Let us just say here that this approach seems to be quite interesting
%the standpoint of system theory.
%Problem~P is tackled in its general form.
%Nevertheless, Cullum is concerned with the approximation of $f^{(1)}$ from
%the samples of $f$.
%Are her results generalizable to the approximation of finitely many
%derivatives of $f$, and if yes, at what expense?
%That is the question!

\section{Illustrative examples.}\setcounter{equation}{0}
\label{examples}

In order to better understand some of the issues involved in the use of  
numerical differentiation and interpolation as a basis for an
observer design method, three academic
examples are currently being pursued.  The first example
involves the estimation
of a time function and a few of its derivatives from noisy samples collected at
discrete instants of time.  The second example builds upon the first one
by showing that if the signal is generated by a known,
continuous-time linear model,
then the model's dynamcis can be incorporated into the estimation
process with advantage.  This will bring us into contact with more classical observer
theory.  The third example will illustrate that the techniques sketched in the
first two examples are
also applicable to nonlinear system models, which is of course, the main
point.


In each case, polynomials are used as the interpolating
funtions, total least squares at the interpolating points is the 
metric to be minimized, and the discrete data will be gathered at a uniform rate.
These simplifications are made so that the discussion can be focused on other issues.

It is emphasized that no theorems are formulated and thus no proofs of validity
are
offered for any of the algorithms proposed herein.  This paper is meant to be an
exploration of a set of ideas.  Formalization of these ideas into a rigorous design
methodology will be pursued when sufficient numerical evidence exists suggesting that
it is useful to do so.  

\subsection{Sum of two sinusoids.}

The purpose of this example is to illustrate the use of numerical
differentiation, alone and 
in combination with standard system theoretic notions,
to recover or estimate several derivatives of a signal in noise.  The key
parameters of interest are: $N$, the order of the interpolating polynomial;
$\Delta t$, the time interval between data points;
$W * \Delta t$, the length in time of the moving window used for data interpolation
(note that $W + 1$ is the number of data points in the window);
$K$, the data node or ``knot'' within the moving window 
where the derivative is to be estimated (counted from the left, with
the first node in the window numbered zero); and $\sigma$, the ``intensity'' 
of the noise process corrupting the measurements.

%The insight developed here will be transferred to an observation
%problem for a dynamical system.

Let an analog signal be given by the sum of two sinusoids of frequencies 1 Hz. and
5 Hz., respectively:
\begin{equation}
y(t) = sin(2 \pi t) + sin(10 \pi t).
\label{sines}
\end{equation}



\begin{figure}[hbt]
\centerline{\psfig{file=sumofsines.eps,width=6in}}
\smallskip
\small
\caption{Time plots of $y(t)$, $y(t) + w(t)$, $y^m_k$ and $y^m_k + w_k$, respectively.}
\label{signals}
\end{figure}


Assume that the measured signal, $y^m$, is the sum of $y$ and $w$, where $w$ is a
Gaussian white noise process, with standard deviation $\sigma$.  
It is further supposed that
$y^m$ is to be sampled at discrete instants of time,
$k \Delta t$, $k = 0, 1, ...$.  In the absence of noise, the minimum sampling rate 
to reconstruct $y$ from sampled data would be the Nyquist rate, 10 Hz.; with
noise,
at least four times this rate is warranted (an anti-aliasing filter is not
being used
though it would be standard practice to do so).
%; its passband will be chosen as 20 Hz.  
%It is then good practice to oversample the filtered signal, by a factor
%of two in order to account for the non-ideal nature of the
%anti-aliasing filter.   
Discrete samples of $y + w$ will be collected therefore
at 40 Hz.:
\begin{eqnarray}
y^m(t) &:=& y(t) + w(t) \\
y^m_k & :=& y^m(k \Delta t). \nonumber
\label{ysampled}
\end{eqnarray}
%where $G_{BW}$ is the Butterworth, anti-aliasing filter.  
The time signals discussed are depicted in Figure \ref{signals}.  The noise
intensity was selected to provide a (numerically computed)
signal to noise ratio of 20/1, and
corresponded to $\sigma = 0.22$. This should 
provide a feel for the sample rate and noise intensity being used in all later
simulations.


Let the interpolating polynomial for the window of 
data $ \{ y_{k-W}^m, \ldots, y_k^m \}$
be denoted by
\begin{equation}
\hat{y}_k(t) = a_0 + a_1 (t - t_{k-W}) + \ldots + a_N (t - t_{k-W})^N,
\label{interppoly}
\end{equation}
where $t_k := k \Delta t$.
The coefficients $\{ a_0, a_1, \ldots, a_N \}$ are determined from the least
squares solution of 
\begin{eqnarray}
\left[ \begin{array}{cccc} 1 & 0 & \ldots & 0 \\
1 & \Delta t & \ldots & (\Delta t)^N \\
\vdots & \vdots & \vdots & \vdots \\
1 & W \Delta t & \ldots & (W \Delta t)^N 
\end{array} \right] \left[ \begin{array}{l} a_0 \\ a_1 \\ \vdots \\ a_N
\end{array} \right] & = & \left[ \begin{array}{l} y^m_{k-W} \\
y^m_{k-W+1} \\ \vdots \\ y^m_k
\end{array} \right],
\label{leastsquares}
\end{eqnarray}
with respect to the Euclidean norm.  
The estimates of the derivatives of $y$ at time $\bar{t}$ are determined by
\begin{equation}
\widehat{ \frac{d^j}{dt^j}y_k(\bar{t}) } := \frac{d^j}{dt^j} \hat{y}_k(t)_{|t=\bar{t}} \  ;
\label{derivestimates}
\end{equation}
for simplicity of notation, this is written as 
$\frac{d^j}{dt^j} \hat{y}_k( \bar{t} )$.


A number of numerical experiments or simulations
were conducted to ascertain the effects of the
parameters $N$,  $\Delta t$, $W$, $K$ and $\sigma$ on the ability to 
estimate the derivatives
of $y$. 
The general trends of these experiments are summarized here; quantifying
these observations analytically would be a worthy goal:
\begin{itemize}

\item The order of the interpolating polynomial should be selected as low as 
possible in order to smooth the noise in the signal. 
On the other hand, for the estimation of 
the 3rd derivative of a signal, for example, 
at least a 3rd order polynomial is required.

\item The window length $W \Delta t$ must be chosen large enough to
capture the variations of the signal so that the interpolating polymomial 
can reproduce its derivatives.  However, larger $W \Delta t$ require
higher order polynomials to be chosen.

\item Estimates of the derivatives made at the end-points of the window
$W \Delta t$ are considerably less accurate than those made in the interior.
However, selecting $K < W$ introduces a time delay of $(W-K) \Delta t$
in the estimates of the signal and its derivatives. 

\item For fixed $N$ and $W \Delta t$, decreasing the time between samples
$\Delta t$ or increasing
the number of points in the window $W$, allows a higher noise intensity
$\sigma$ to be tolerated.

\item For a fixed $\Delta t$, though increasing $W$ can produce more accurate results,
it must be considered that it requires more on line compuation and a larger
delay must be introduced (i.e., $(W-K)\Delta t$ increased)
in order to retain the accuracy of the estimates. 

\end{itemize}

\begin{figure}[hbt]
\centerline{\psfig{file=mplot11.eps,width=6in}}
\smallskip
\small
\caption{Comparison of the true and estimated $(\times)$ values of
$y, \ \frac{dy}{dt}, \
\frac{d^2y}{dt^2}, \mbox{and} \
\frac{d^3y}{dt^3}$, respectively using interpolation/numerical differentiation.}
\label{figmplot10}
\end{figure}

Figure \ref{figmplot10} compares the true and estimated values of
$y, \ \frac{dy}{dt}, \
\frac{d^2y}{dt^2}, \mbox{and} \
\frac{d^3y}{dt^3}$ when $N=4$, $\Delta t = .025$, $W=8$ and $K=5$. 
Figure \ref{figsineinterp} displays the normalized errors
in the estimates: the errors in the estimates 
for $(y, \ldots, \frac{d^3}{dt^3}y)$ are normalized
by 
\begin{equation}
lim_{T \rightarrow \infty} \frac{1}{T} \int_{0}^{T}{|\frac{d^j}{dt^j}y(t)| dt}.
\label{scales}
\end{equation}
Thus, the error in $y$ is divided by 0.796, $\frac{dy}{dt}$ by 19.3,
$\frac{d^2y}{dt^2}$ by 650, 
and $\frac{d^3y}{dt^3}$ by 22,700.  The wide range in the magnitudes makes the
estimation problem quite challenging.


\begin{figure}[hbt]
\centerline{\psfig{file=mplot9.eps,width=6in}}
%\centerline{\psfig{file=sine_interp.eps}}
\smallskip
\small
\caption{Normalized estimation errors in $y, \ \frac{dy}{dt}, \
\frac{d^2y}{dt^2}, \mbox{and} \
\frac{d^3y}{dt^3}$, respectively using interpolation/numerical differentiation.}
\label{figsineinterp}
\end{figure}


Estimating the signal and its derivatives on the basis of 
the $W + 1$ data points in the moving window corresponds to using
an FIR (finite impulse response) 
filter.  By introducing a simple modification
to the above, a sort of IIR (infinite
impulse response) filter can be achieved.  The modification is somewhat
analogous to 
what is done in spline
approximations, but to a systems person, it would be called a state.

The idea is to append to (\ref{leastsquares}), the estimates of $y$ and
its derivatives at node K from the previous window, that is, one appends
the relations
\begin{equation}
\frac{d^j}{dt^j} \hat{y}_k(t_{k-1-W+K}) =  \frac{d^j}{dt^j} \hat{y}_{k-1}(t_{k-1-W+K}),
\label{derivconstraints}
\end{equation}
for $j=0, \ldots, J$, for some $0 \le J \le N$.
This is equivalent to
\begin{eqnarray} {\tiny
\left[ \begin{array}{ccccccc} 
1 & (K-1) \Delta t & ((K-1)\Delta t)^2 & \ldots & ((K-1) \Delta t)^J & 
\ldots & ((K-1) \Delta t)^N \\
0 & 1 & 2(K-1) \Delta t & \ldots & J((K-1) \Delta t)^{J-1} & \ldots & 
N((K-1) \Delta t)^{N-1} \\
\vdots & \vdots & \vdots & \vdots & \vdots & \vdots & \vdots \\
0 & 0 & \ldots & \ldots & J! & \ldots & N(N-1)\ldots(N-J+1)((K-1) \Delta t)^{N-J}
\end{array} \right] } \times & \mbox{ } &  \nonumber \\
\left[ \begin{array}{l} a_0 \\ a_1 \\ \vdots \\ a_N \end{array} \right] 
 =  \left[ \begin{array}{r} 
\hat{y}_{k-1}(t_{k-1-W+K}) \\
\frac{d}{dt} \hat{y}_{k-1}(t_{k-1-W+K}) \\ 
\vdots \\ 
\frac{d^J}{dt^J} \hat{y}_{k-1}(t_{k-1-W+K})
\end{array} \right] & \mbox{\hspace*{.2in}} &.
\label{matrixderivconstraints}
\end{eqnarray}

The right-hand side of (\ref{derivconstraints}) must be initialized to
start the interpolation process, and thus becomes a state.  This state 
allows derivative 
information to be passed from one window of data 
to the next.   
%Often, a smaller window can be used to achieve similar noise attenuation.  
Weights can be added to trade-off errors in interpolating the measured data points
versus errors in interpolating the derivatives in (\ref{derivconstraints}).  



\subsection{Including a system model in the interpolation process.}

The goal of this example is to demonstrate two methods for incorporating
a dynamic model of a system into the interpolation/numerical differentiation
process,
whenever the signal in question is generated by a finite set of ordinary
differential equations.
This will bring us into more direct contact with observer design because we
will tightly connect estimating the derivatives of a system's output
with estimating the states of the model.

The signal (\ref{sines}) can obviously be generated by an initialized,
uncontrolled, linear model with four states.  Let 
\begin{equation}
\begin{array}{rcl}
\frac{dx}{dt} & = &  A x  \\
y & = & C x 
\end{array}
\label{linearmodel}
\end{equation}
be such a model; it will be observable.  This means that $x$ can be recovered
from estimates of $y$ and its derivatives through
\begin{equation} \left[ \begin{array}{r} 
y \\ \frac{d}{dt}y \\ \vdots \\ \frac{d^{\eta}}{dt^{\eta}}y
\end{array} \right] = \left[ \begin{array}{l} 
C \\ CA \\ \vdots \\ CA^{\eta}
\end{array} \right] x
\label{obsmat}
\end{equation}
for $\eta$ large enough, which in the particular case of (\ref{linearmodel}) is
$\eta = 3$.

The first way to include the model in the estimation process is now
explained.  Suppose that in general the model (\ref{linearmodel}) has dimension
$n$, and for simplicity, assume that it has a single output.
As before, let
\begin{equation}
\left[ \begin{array}{r} 
\hat{y}_{k-1}(t_{k-1-W+K}) \\
\frac{d}{dt} \hat{y}_{k-1}(t_{k-1-W+K}) \\ 
\vdots \\ 
\frac{d^{n-1}}{dt^{n-1}} \hat{y}_{k-1}(t_{k-1-W+K})
\end{array} \right]
\label{ydotest}
\end{equation}
be estimates of $y$ and its first $n-1$ derivatives from the previous window.
By substituting (\ref{ydotest}) for the right-hand side of (\ref{obsmat}) with
$\eta = n-1$, one can
determine $\hat{x}_{k-1}(t_{k-1-W+K})$. 
If the resulting estimate of $x$ is ``accurate'', then it satisfies (\ref{obsmat}) for
any $\eta$; if not, choosing $\eta \ge n$ introduces constraints whose
errors can be minimized in order for $x$ to be more closely compatible
with the dynamics (\ref{linearmodel}).  Thus, in place of
(\ref{derivconstraints}),
one augments (\ref{leastsquares}) with
\begin{equation}
\frac{d^j}{dt^j} \hat{y}_k(t_{k-1-W+K}) = CA^j  \hat{x}_{k-1}(t_{k-1-W+K}),
\label{modelconstraints}
\end{equation}
For $j=0, \ldots, \eta$, and proceeds as indicated previously.  
Note that taking $\eta =J \le n-1$ reduces 
(\ref{modelconstraints}) to (\ref{derivconstraints}).



\begin{figure}[hbt]
\centerline{\psfig{file=sine_spline_order5.eps,width=6in}}
\smallskip
\caption{ Normalized estimation errors in $y, \ \frac{dy}{dt}, \
\frac{d^2y}{dt^2}, \mbox{and} \
\frac{d^3y}{dt^3}$, respectively using interpolation/numerical differentiation
and $\eta=5$ in (\protect\ref{modelconstraints}).  Note the reduced error with respect to
Figure \protect\ref{figsineinterp}. }
\label{figsinesplineorder5}
\end{figure}


Figure \ref{figsinesplineorder5} shows the
effect of this modification on the estimation process for $N=6$, $\Delta t = .025$, 
$W=8$ and $K=5$.  The weight on the constraints was chosen as 
\begin{equation} Q = 
\left[ \begin{array}{cccccc} 1.26 & 0 & 0 & 0 & 0 & 0\\
0  & 5.19 \ 10^{-2}  & 0 & 0 & 0 & 0 \\
0 & 0 & 1.54 \ 10^{-3} & 0 & 0 & 0 \\
0 & 0 & 0 & 4.4 \ 10^{-5} & 0 & 0 \\
0 & 0 & 0 & 0 & 1.3 \ 10^{-6} & 0 \\
0 & 0 & 0 & 0 & 0 & 4.4 \ 10^{-8}
\end{array} \right].
\end{equation}
The rationale for
selecting the entries of the
weight $Q$ is that the magnitude of each successive derivative grows by a
approximatley
30 (recall (\ref{scales}), or see (\ref{sines})), 
and thus $Q$ corresponds to equal ``relative
weighting'' on $y$ and its derivatives;
this same rationale will be applied to subsequent examples. Including the model
interpolation constraints allowed the same window size and delay as in Figure
\ref{figsineinterp} to be used with an increased polynomial order, without
``following'' the noise.


The method just outlined for estimating $x$ from $y$ and its derivatives is akin to
a ``dead-beat'' observer for $x$, though, as explained
earlier, there is smoothing in the
interpolation of $y$, which is ``IIR''.  Some additional smoothing from the
model can be obtained, and this constitutes the second method for introducing
the model into the interpolation/numerical differentiation process.
This technique, which is reminiscent of
\cite{grizzle1990b,grizzle1990a,grizzle1990c,moraal1992}, 
will be
directly applicable to nonlinear systems.  

Let $A_d$ denote the discretization of the linear model, (\ref{linearmodel}),
and let $\hat{x}_{k-1}^{+}(t_{k-1-W+K})$ be the previous estimate of $x$.  Update $x$
based on the latest estimate of $y$ and its derivatives through a damped Newton
method:
\begin{eqnarray}
\hat{x}_{k}^{-}(t_{k-W+K}) & = & A_d \hat{x}_{k-1}^{+}(t_{k-1-W+K}) \nonumber \\
\hat{x}_{k}^{+}(t_{k-W+K}) & = & \hat{x}_{k}^{-}(t_{k-W+K}) + \label{newton} \\
\epsilon \left[ \begin{array}{l} C \\ CA \\ \vdots \\ CA^{n-1} \end{array} 
\right]^{-1} & &  ( \left[ \begin{array}{l} 
\hat{y}_{k}(t_{k-W+K}) \\
\frac{d}{dt} \hat{y}_{k}(t_{k-W+K}) \\ 
\vdots \\ 
\frac{d^{n-1}}{dt^{n-1}} \hat{y}_{k}(t_{k-W+K}) \end{array} \right] - 
\left[ \begin{array}{l} C \\ CA \\ \vdots \\ CA^{n-1} \end{array} 
\right]  \hat{x}_{k}^{-}(t_{k-W+K}) ), \nonumber
\end{eqnarray}
where $0 < \epsilon <  1$; using $\epsilon = 1$ reduces the above to a ``dead-beat''
estimator for $x$.  The damped Newton update can be used with
(\ref{leastsquares}) alone or with any of
the additions to (\ref{leastsquares}) discussed previously.

Selected simulations of the above estimation methods are now presented.
In what follows,
the state in (\ref{linearmodel}) was chosen as 
$$ x = \left[ \begin{array}{r}
y \\ \frac{d}{dt} y \\  \frac{d^2}{dt^2} y \\ \frac{d^3}{dt^3} y 
\end{array} \right]. $$  
It was properly initialized so that its output is precisely
given by (\ref{sines}).  The same noise sequence used in
Figure \ref{figsineinterp} was added to the output of (\ref{linearmodel}).
%also, the anti-aliasing filter (\ref{ysampled}) was used.


\begin{figure}[hbt]
\centerline{\psfig{file=sine_spline.eps,width=6in}}

\caption{Normalized estimation errors in \protect$y, \ \frac{dy}{dt}, \
\frac{d^2y}{dt^2}$ 
and  \protect$\frac{d^3y}{dt^3}$, respectively, incorporating 
``spline method'' and
Newton update in interpolation/numerical differentiation.}
\label{figsinespline}
\end{figure}


Figure \ref{figsinespline} shows the normalized estimation error resulting
from the application of (\ref{leastsquares})  in combination with
(\ref{derivconstraints})
and (\ref{newton}).   The parameter values were chosen as: $N=5$, $\Delta t = 0.025$
sec., $W=6$ and $K=3$  in (\ref{leastsquares}), $\eta = 4$ in 
(\ref{modelconstraints}) with
a weight on (\ref{derivconstraints}) selected as
\begin{equation} Q = 0.1 \times
\left[ \begin{array}{cccc} 1.26 \, 10^{-1} & 0 & 0 & 0 \\
0  & 5.19 \, 10^{-3}  & 0 & 0 \\
0 & 0 & 1.54 \, 10^{-4} & 0 \\
0 & 0 & 0 & 4.4 \, 10^{-5} 
\end{array} \right],
\end{equation}
and $\epsilon = 0.03$ in (\ref{newton}).   
For comparison purposes,
a steady state, discrete-time, Kalman filter was designed with 
state noise covariance equal to 
$0.01 \ \sigma^2 \ Q^{-1}$ and output noise variance equal to $\sigma^2$; even
though (\ref{linearmodel}) does not have any process noise per se, the pair
consisting of the matrix $A$ and the state noise convariance matrix
had to be made stabilizable
for a stable Kalman filter to exist.
The magnitudes of the steady state errors were essentially identical to those in 
Figure \ref{figsinespline}.


\begin{figure}[hbt]
\centerline{\psfig{file=sine_model.eps,width=6in}}
%\centerline{\psfig{file=sine_model.eps}}

\caption{Normalized estimation errors in $y, \ \frac{dy}{dt}, \
\frac{d^2y}{dt^2} 
\mbox{and} \ \frac{d^3y}{dt^3}$, respectively, incorporating model derivative
constraints and damped Newton updates with 
interpolation/numerical differentiation.}
\label{figsinemodel}
\end{figure}


The normalized estimation errors resulting
from the application of (\ref{leastsquares}) in combination with
(\ref{modelconstraints})
and (\ref{newton}) are shown in
Figure \ref{figsinemodel}.   The parameter values were chosen as: $N=5$, 
$\Delta t = 0.025$
sec., $W=6$ and $K=3$  in (\ref{leastsquares}), $\eta = 4$ in
(\ref{modelconstraints}) 
with a weight selected as
\begin{equation} Q = 
\left[ \begin{array}{ccccc} 1.26 \, 10^{-1} & 0 & 0 & 0 & 0 \\
0  & 5.19 \, 10^{-3}  & 0 & 0 & 0 \\
0 & 0 & 1.54 \, 10^{-4} & 0 & 0 \\
0 & 0 & 0 & 4.4 \, 10^{-5} & 0 \\
0 & 0 & 0 & 0 & 1.33 \, 10^{-6}
\end{array} \right];
\end{equation}
$\epsilon$ in (\ref{newton}) was set at $0.03$.  

\subsection{Third order, nonlinear dynamical system.}

The goal of this example is to show that the method illustrated on the
linear example above is also applicable to nonlinear models.  For this
purpose, consider the following system which is a ``convex combination''
of an unstable linear dynamics and a stable linear dynamics:
\begin{eqnarray}
\left[ \begin{array}{c} \frac{dx_1}{dt} \\ \frac{dx_2}{dt}  
\\ \frac{dx_3}{dt} \end{array} \right] 
& = & f(x) = 
\left[ \begin{array}{c} x_2 \\ x_3 \\ f_3(x)
\end{array} \right]   \label{nonlinmodel} \\
y & = & h(x) = x_1 , \nonumber 
\label{dynamic}
\end{eqnarray}
where,
\begin{eqnarray*}
f_3(x) & = & \alpha(x) g_s(x) + (1 - \alpha(x)) g_u(x), \\
\alpha(x)& =& \beta  \frac{x'x}{1 + x'x} \\ 
g_s(x) & = & -54 x_1 - 36 x_2 - 9 x_3 \\
g_u(x) & = & 54 x_1 - 36 x_2 + 9 x_3.
\end{eqnarray*}
The parameter $\beta$ will be fixed at 2 for now; in the final version of the
paper, it will be assumed unknown.  

\begin{figure}[thb]
\centerline{\psfig{file=dyn_sampled.eps,width=6in}}
\smallskip
\caption{Time plots of $y(t)$, $y(t) + w(t)$, $y^m_k$ and $y^m_k + w_k$, 
respectively.}
\label{figdynamicsampled}
\end{figure}

The output of system (\ref{dynamic}) has a fundamental frequency 
of about 1 Hz., so the sample rate for the system was selected as
4 Hz.; this is probably a little too slow, but will
challenge the observers to be designed.
Figure \ref{figdynamicsampled} shows the output of the system (\ref{dynamic}),
$y$, $y + w$, $y^m_k$ and $y^m_k + w_k$, respectively, for $\Delta t = 0.25$
sec. and  $\sigma= 0.0168$ (which
corresponds once again to a 20/1 signal to noise ratio). Figure
\ref{figdynamicinterp} shows the effectiveness of estimating the states
on the basis of (\ref{leastsquares}) with $N = 6$,
$W=8$, $K=4$.  The performance seems quite good.

\begin{figure}[hbt]
\centerline{\psfig{file=mplot12.eps,width=6in}}
\smallskip
\caption{Comparison of estimated signals $(\times)$'s and true
signals using interpolation/numerical differentiation.}
\label{figdynamicinterp}
\end{figure}

A state estimator was also constructed using the interpolation/numerical
differentiation method of (\ref{leastsquares}), with  (\ref{newton}) replaced by 
\begin{eqnarray}
\hat{x}_{k}^{-}(t_{k-W+K}) & = & F_{\Delta t}(\hat{x}_{k-1}^{+}(t_{k-1-W+K}))
\nonumber \\
\hat{x}_{k}^{+}(t_{k-W+K}) & = & \hat{x}_{k}^{-}(t_{k-W+K})+ \label{nonlinnewton} \\ 
\epsilon \, (\frac{\partial}{\partial x}
\left[ \begin{array}{l} h \\ L_fh \\ L_f^2h \end{array} 
\right]) ^{-1}_{|\hat{x}_{k}^{-}(t_{k-W+K})} & & ( \left[ \begin{array}{r} 
\hat{y}_{k}(t_{k-W+K}) \\
\frac{d}{dt} \hat{y}_{k}(t_{k-W+K}) \\ 
\frac{d^{2}}{dt^{2}} \hat{y}_{k}(t_{k-W+K}) \end{array} \right] - 
\left[ \begin{array}{l} h \\ L_fh \\ L_f^2h \end{array} 
\right]_{| \hat{x}_{k}^{-}(t_{k-W+K})} ). \nonumber
\end{eqnarray}
In the above, $F_{\Delta t}(x)$ is the 
sampled-data representation of (\ref{nonlinmodel}) for $t=\Delta t$; for the
estimator results which follow, it was computed by Euler
integration with step size $0.0125$ (i.e., 80 Hz. updates).
The parameter $\epsilon$ was selected as
$$ \epsilon = \left[ \begin{array}{ccc}
0.6 & 0 & \\
0 & 0.12 & 0 \\
0 & 0 & 0.03  \end{array} \right], $$ to reflect scaling in the
magnitudes of the estimates of the derivatives and uncertainty due to noise.
Figure \ref{figdynamicnewton} compares the estimated values
of $x_1, x_2$ and $x_3$ to the true values.



\begin{figure}[hbt]
\centerline{\psfig{file=dynamic_newton.eps,width=6in}}
\smallskip
\caption{Comparison of estimated signals $(\times)$'s and true
signals using (\protect\ref{nonlinnewton}).  Note the reduced error with
respect to Figure \protect\ref{figdynamicinterp}. }
\label{figdynamicnewton}
\end{figure}


One could explore replacing the model 
constraints (\ref{modelconstraints}) by
\begin{equation}
\frac{d^j}{dt^j} \hat{y}_k(t_{k-1-W+K}) = L_f^j h(\hat{x}_{k-1}(t_{k-1-W+K})),
\label{nonlinmodelconstraints}
\end{equation}
for $j = 0, \ldots, J$, but this was not done.


%\section{Conclusions}\setcounter{equation}{0}
%\label{conclusions}

\bibliographystyle{plain}


\bibliography{numdiff}


\end{document}



